\section{The cubic digamma finite part and the first harmonic coefficient}
\label{bridge:sec}

In this section write $\eta=\eta(1)$ and $Z_2(s,1)=Z_2(s,1;1)$,
with the conventions fixed in Section~\ref{sec:scope}. In particular,
$B_2=(\gamma^2-\zeta(2))/2$ and
\[
 Z_2(1+u,1)=u^{-2}+\gamma u^{-1}+B_2+\eta u+O(u^2).
\]
The inclusive harmonic coefficient is related by
$\eta=-\gamma_H(1,1)-\zeta'(2)$ \cite{Kargin,IncomingOrdered}.

\begin{theorem}[Cubic--harmonic bridge]\label{bridge:main}
In the unit-coordinate Hadamard finite-part convention,
\begin{align}
 Q:=\operatorname{FP}\int_0^1\psi(x)^3\,dx
 &=3\zeta(2)-6\eta-6\gamma\gamma_1,\label{bridge:eta}\\
 &=3\zeta(2)+6\gamma_H(1,1)+6\zeta'(2)-6\gamma\gamma_1.
 \label{bridge:harmonic}
\end{align}
An absolutely convergent arithmetic representation is
\begin{equation}\label{bridge:arithmetic}
 Q=3\zeta(2)+6\gamma_2+
  6\sum_{n=1}^{\infty}
   \frac{\bigl(H_{n-1}-\gamma-\log n\bigr)\log n}{n}.
\end{equation}
Consequently the cubic finite part and the first strict harmonic
coefficient are equivalent modulo ordinary zeta and Stieltjes constants.
This theorem does not assert that the harmonic coefficient is a finite
combination of ordinary zeta constants.
\end{theorem}

\subsection{A global partial-fraction identity}

For $n\ge0$ put
\[
 b_n=H_n-\gamma,\qquad
 a_n=\zeta(2)+H_n^{(2)}-b_n^2,
 \qquad H_n^{(r)}=\sum_{j=1}^n j^{-r},
\]
with $H_0^{(r)}=0$.  Set
\[
 c_0=-\gamma^3+6\gamma\zeta(2)-3\zeta(3).
\]

\begin{lemma}[Normally convergent cubic partial fractions]
\label{bridge:ML}
For $z\notin\{0,-1,-2,\ldots\}$,
\begin{align}
 \psi(z)^3={}&-\frac1{z^3}-\frac{3\gamma}{z^2}
       +\frac{3(\zeta(2)-\gamma^2)}z+c_0\notag\\
 &+\sum_{n=1}^{\infty}\bigg\{
      -\bigl((z+n)^{-3}-n^{-3}\bigr)
      +3b_n\bigl((z+n)^{-2}-n^{-2}\bigr)\notag\\
 &\hspace{42mm}+3a_n\bigl((z+n)^{-1}-n^{-1}\bigr)\bigg\}.
 \label{bridge:MLformula}
\end{align}
The series converges normally on every compact subset avoiding the poles.
\end{lemma}

\begin{proof}
At $z=-n+\varepsilon$, the recurrence of $\psi$ gives
\[
 \psi(-n+\varepsilon)
   =-\varepsilon^{-1}+b_n+
      \bigl(\zeta(2)+H_n^{(2)}\bigr)\varepsilon
      +\bigl(H_n^{(3)}-\zeta(3)\bigr)\varepsilon^2
      +O(\varepsilon^3).
\]
The cubic principal part at $-n$ is therefore
\[
 P_n(z)=-(z+n)^{-3}+3b_n(z+n)^{-2}+3a_n(z+n)^{-1}.
\]
The constant coefficient at zero is $c_0$.
It is important to justify the entire remainder in a Mittag--Leffler
argument, rather than infer it from the principal parts alone.
Let $f(z)=\psi(z)^3-P_0(z)$, which is regular at zero, and integrate
\[
 \frac{z f(w)}{w(w-z)}
\]
over the circles $|w|=N+\tfrac12$, for fixed $z$ and large $N$.
On these circles $f(w)=O(\log^3 N)$ uniformly.  In the right
half-plane this follows from the usual digamma asymptotic \cite[Sections~5.7, 5.11]{DLMFgamma}.  In the
left half-plane use
\[
 \psi(w)=\psi(1-w)-\pi\cot(\pi w).
\]
The cotangent is uniformly bounded: when $|\Im w|\ge1/4$ this follows
from its elementary exponential formula; in the remaining part of the
circle, the real coordinate stays a fixed positive distance from the
integers.  Thus the contour integral is
$O(\log^3 N/N)$ and tends to zero.
Its residues at $z$ and $0$ are $f(z)$ and $-f(0)$.
The residue at $-n$ is $-\bigl(P_n(z)-P_n(0)\bigr)$, either by a
direct calculation or by applying the same two-point Cauchy kernel to
each principal part.  The residue theorem proves
$f(z)-f(0)=\sum_{n\ge1}(P_n(z)-P_n(0))$.
Finally $b_n=O(\log n)$ and $a_n=O(\log^2 n)$, while on compact sets
$(z+n)^{-1}-n^{-1}=O(n^{-2})$.  These estimates prove normal
convergence and justify the asserted identity.
\end{proof}

\subsection{Reduction to one logarithmic harmonic series}

Integrating the normally convergent remainder of
\eqref{bridge:MLformula} term by term yields
\begin{align}
 Q={}&\frac12+3\gamma+c_0
  +\sum_{n\ge1}\bigg\{
    \frac1{n^3}+\frac1{2(n+1)^2}-\frac1{2n^2}
     -\frac{3b_n}{n^2(n+1)}\notag\\
 &\hspace{47mm}
     +3a_n\left(\log\left(1+\frac1n\right)-\frac1n\right)
     \bigg\}.
 \label{bridge:integratedML}
\end{align}
Here the finite parts of $-x^{-3}$ and $-3\gamma x^{-2}$ are
$1/2$ and $3\gamma$, and the unit-coordinate finite part of
$x^{-1}$ is zero.  The elementary identities
\[
 \sum_{n\ge1}\frac{H_n}{n(n+1)}=\zeta(2),\qquad
 \sum_{n\ge1}\frac{H_n}{n^2}=2\zeta(3)
\]
give
\[
 \sum_{n\ge1}\frac{b_n}{n^2(n+1)}
   =2\zeta(3)-\zeta(2)-\gamma\bigl(\zeta(2)-1\bigr).
\]
For completeness, the second Euler sum follows from
\[
 \sum_{n\ge1}\frac{H_n}{n^2}
 =\int_0^1\frac{\log x\log(1-x)}{x(1-x)}\,dx
 =2\int_0^1\frac{\log x\log(1-x)}x\,dx
 =2\zeta(3),
\]
using the harmonic generating function, symmetry, and an absolutely
convergent geometric-series integration.
Consequently \eqref{bridge:integratedML} becomes
\begin{equation}\label{bridge:quadratic-log}
 Q=-\gamma^3+9\gamma\zeta(2)+3\zeta(2)-8\zeta(3)
  +3\sum_{n\ge1}a_n
       \left(\log\left(1+\frac1n\right)-\frac1n\right).
\end{equation}

The crucial cancellation is the exact difference
\[
 a_n-a_{n-1}=-\frac{2(H_{n-1}-\gamma)}n.
\]
Write
$L_N=\sum_{n=1}^N(H_{n-1}-\gamma)\log n/n$ and
$S_N=\sum_{n=1}^N H_n/n^2$.  Finite summation by parts gives
\begin{align*}
 \sum_{n=1}^N a_n\log\left(1+\frac1n\right)
 &=a_N\log(N+1)+2L_N,\\
 \sum_{n=1}^N\frac{a_n}{n}
 &=(\zeta(2)-\gamma^2)H_N+
      \gamma(H_N^2+H_N^{(2)})-\frac{H_N^3}{3}
       +H_NH_N^{(2)}-2S_N+\frac43H_N^{(3)}.
\end{align*}
The second identity follows by summing the finite differences of
$H_n^2$, $H_n^3$, and $H_nH_n^{(2)}$; it has no limiting
interchange hidden in it.  Since
\[
 H_N-\gamma=\log(N+1)+O(N^{-1}),\quad
 H_N^{(2)}=\zeta(2)+O(N^{-1}),\quad S_N\longrightarrow2\zeta(3),
\]
substitution in \eqref{bridge:quadratic-log}, with all polynomial
terms retained before the limit, gives
\begin{equation}\label{bridge:limit}
 Q=3\zeta(2)+6\lim_{N\to\infty}
       \left(L_N-\frac{\log^3 N}{3}\right).
\end{equation}
Now split $H_{n-1}-\gamma=\log n+\psi(n)-\log n$ and use the
standard defining limit
\[
 \gamma_2=\lim_{N\to\infty}
   \left(\sum_{n=1}^N\frac{\log^2 n}{n}-\frac{\log^3 N}{3}\right).
\]
Since $\psi(n)-\log n=O(n^{-1})$, the remaining series is absolutely
convergent.  This proves \eqref{bridge:arithmetic}.

Define the holomorphic Dirichlet deformation
\[
 D(s)=\sum_{n\ge1}\frac{\psi(n)-\log n}{n^s},\qquad\Re s>0.
\]
Normal convergence, including every fixed derivative on compact
subsets, follows from $\psi(n)-\log n=O(n^{-1})$.
In $\Re s>1$ we have the exact identity
\[
 D(s)=Z_2(s,1)-\gamma\zeta(s)+\zeta'(s).
\]
Comparing Laurent coefficients at $s=1$ gives
$D'(1)=\eta+\gamma\gamma_1+\gamma_2$.  Formula
\eqref{bridge:arithmetic} says
$Q=3\zeta(2)+6\gamma_2-6D'(1)$, proving
\eqref{bridge:eta}.  The inclusive-to-strict conversion gives
\eqref{bridge:harmonic} and completes the proof of
Theorem~\ref{bridge:main}.

\subsection{A normalized complex antiderivative}

\begin{corollary}[Cubic digamma primitive]\label{bridge:primitive}
On $\Re z>0$, with the principal logarithm, the function
\begin{align}
 \mathcal P_3(z)={}&\frac1{2z^2}+\frac{3\gamma}z
       +3(\zeta(2)-\gamma^2)\log z+c_0z\notag\\
 &+\sum_{n=1}^{\infty}\bigg\{
      \frac1{2(n+z)^2}-\frac1{2n^2}+\frac{z}{n^3}
      +3b_n\left(\frac1n-\frac1{n+z}-\frac z{n^2}\right)
        \notag\\
 &\hspace{41mm}
      +3a_n\left(\log\left(1+\frac zn\right)-\frac zn\right)
        \bigg\}
 \label{bridge:primitiveformula}
\end{align}
satisfies $\mathcal P_3'(z)=\psi(z)^3$.  Its asymptotic expansion at
zero has constant term zero.  In particular, for every real $b>0$,
\[
 \operatorname{FP}\int_0^b\psi(x)^3\,dx=\mathcal P_3(b),
 \qquad \mathcal P_3(1)=3\zeta(2)-6\eta-6\gamma\gamma_1.
\]
The series and its differentiated series converge normally on compact
subsets of the right half-plane.  Higher derivatives therefore give
identities for derivatives of $\psi^3$ in the same domain.
\end{corollary}

\begin{proof}
Each brace is the primitive, vanishing at $z=0$, of the corresponding
brace in \eqref{bridge:MLformula}.  On a fixed compact set,
the logarithmic difference is $O(n^{-2})$, the quadratic-pole
primitive is $O(n^{-3})$, and the cubic-pole primitive is $O(n^{-4})$.
Multiplication by $a_n=O(\log^2 n)$ or $b_n=O(\log n)$ preserves
normal convergence.  Differentiation recovers the previous lemma.
All series terms vanish at zero, so the displayed singular and logarithmic
terms leave no constant term.  This proves the stated normalization.
\end{proof}

\begin{corollary}[A polygamma--harmonic hierarchy]\label{bridge:polygamma}
For every integer $r\ge1$ and $z\notin\{0,-1,-2,\ldots\}$,
\begin{align}
 &\sum_{r_1+r_2+r_3=r}\frac{r!}{r_1!r_2!r_3!}
       \psi^{(r_1)}(z)\psi^{(r_2)}(z)\psi^{(r_3)}(z)\notag\\
 &\quad=(-1)^r r!\sum_{n=0}^{\infty}
   \left\{-\frac{(r+1)(r+2)}{2(n+z)^{r+3}}
       +\frac{3(r+1)b_n}{(n+z)^{r+2}}
       +\frac{3a_n}{(n+z)^{r+1}}\right\}.
 \label{bridge:allpolygamma}
\end{align}
Here $\psi^{(0)}=\psi$. The series converges normally away from
the poles. At $z=1$ the left side is a finite polynomial in
$\gamma,\zeta(2),\ldots,\zeta(r+1)$, using
$\psi^{(k)}(1)=(-1)^{k+1}k!\zeta(k+1)$ for $k\ge1$.
\end{corollary}
\begin{proof}
Differentiate \eqref{bridge:MLformula} $r$ times. The constant
subtractions disappear and the derivative series is bounded on each
compact set by a constant times
$\sum n^{-r-1}(1+\log^2n)$. The product rule gives the left
side. The $n=0$ summand includes the principal part at zero since
$b_0=-\gamma$ and $a_0=\zeta(2)-\gamma^2$. The special values
follow from the classical Hurwitz--polygamma identity
\cite[Section~25.11]{DLMFhurwitz}.
\end{proof}

\subsection{A short certified evaluation route}

For integers $N\ge2$ and $K\ge0$ define
\begin{align*}
 S_{N,K}={}&\sum_{n=1}^{N-1}
  \frac{(H_{n-1}-\gamma-\log n)\log n}{n}
   +\frac12\zeta'(2,N)
   +\sum_{k=1}^K\frac{B_{2k}}{2k}\zeta'(2k+1,N).
\end{align*}
The Binet formula for $\psi$ \cite[Section~5.9]{DLMFgamma} and the finite geometric expansion of
$(n^2+t^2)^{-1}$ give the rigorous truncation estimate
\begin{equation}\label{bridge:tail-bound}
 \left|S_{N,K}-\sum_{n\ge1}
  \frac{(\psi(n)-\log n)\log n}{n}\right|
 \le \frac{|B_{2K+2}|}{2K+2}\bigl(-\zeta'(2K+3,N)\bigr).
\end{equation}
Indeed the digamma asymptotic remainder is bounded by the magnitude of
its first omitted Bernoulli term for positive real arguments, and the
logarithmic tail is exactly a Hurwitz-zeta derivative.
The supplied Arb script evaluates every displayed quantity with ball
arithmetic and adds this analytic tail bound.  This certifies the
arithmetic identity numerically without replacing its proof.

\paragraph{Status.}
The harmonic zeta function, its Stieltjes coefficients, and the
Mittag--Leffler method are classical.  The cubic--harmonic bridge and
its coordinate consequence were derived here and were not found in the
inspected repository reports.  No worldwide priority claim is made.
